% ── Chapter 1: Introduction ───────────────────────────────────
\chapter{บทนำ (Introduction)}

\section{ภูมิหลังและที่มาของปัญหา (Background and Problem Statement)}

การดื้อยาต้านจุลชีพ (Antimicrobial resistance หรือ AMR) เป็นหนึ่งในภัยคุกคามที่เร่งด่วนและซับซ้อนที่สุดต่อระบบสาธารณสุขระดับโลกในศตวรรษที่ 21 เมื่อแบคทีเรีย ไวรัส เชื้อรา หรือพยาธิ เกิดการปรับตัวและพัฒนากลไกต้านทานต่อยาที่ออกแบบมาเพื่อกำจัดเชื้อก่อโรคเหล่านี้ ส่งผลให้การรักษาด้วยวิธีมาตรฐานไม่ได้ผล การติดเชื้อจึงแพร่กระจายได้ง่ายขึ้นและอัตราการเสียชีวิตสูงขึ้น องค์การอนามัยโลก (WHO) ได้จัดให้ AMR เป็นวาระสำคัญเร่งด่วน โดยประเมินว่าการติดเชื้อดื้อยาเป็นสาเหตุโดยตรงที่ทำให้ประชากรราว 1.27 ล้านคนเสียชีวิตในปี 2019 และเป็นสาเหตุทางอ้อมอีกเกือบ 5 ล้านคนทั่วโลก \cite{who_amr2022} หากไม่มีมาตรการรับมือที่บูรณาการร่วมกัน คาดว่า AMR อาจส่งผลให้มีผู้เสียชีวิตสูงถึง 10 ล้านคนต่อปีภายในปี 2050 ซึ่งจะแซงหน้าโรคมะเร็งในฐานะสาเหตุหลักของการเสียชีวิต นอกจากความสูญเสียทางชีวิตแล้ว AMR ยังสร้างภาระทางเศรษฐกิจอย่างมหาศาลจากการที่ผู้ป่วยต้องพักรักษาตัวนานขึ้น ต้องพึ่งพายาทางเลือกที่สองหรือสามที่มีราคาแพง และส่งผลให้ผลิตภาพแรงงานลดลง

กุญแจสำคัญในการรับมือกับ AMR คือการทดสอบความไวต่อยาต้านจุลชีพ (Antimicrobial Susceptibility Testing หรือ AST) ก่อนเริ่มให้การรักษาด้วยยาปฏิชีวนะ แพทย์จำเป็นต้องทราบว่ายาชนิดใดที่ยังมีประสิทธิภาพต่อเชื้อก่อโรคที่พบในผู้ป่วยแต่ละราย การสั่งจ่ายยาที่เชื้อแบคทีเรียดื้ออยู่แล้วไม่เพียงแต่ทำให้การรักษาล้มเหลว แต่ยังเป็นการกระตุ้นให้เกิดเชื้อดื้อยาเพิ่มขึ้นในวงกว้าง ดังนั้น การทำ AST ที่แม่นยำและทราบผลอย่างทันท่วงทีจึงไม่ใช่เพียงขั้นตอนการวินิจฉัยโรค แต่ยังเป็นมาตรการทางสาธารณสุขที่สำคัญ ความล่าช้าในการรอผลทดสอบ ซึ่งปกติกินเวลา 24 ถึง 72 ชั่วโมงด้วยวิธีเพาะเชื้อดั้งเดิม บังคับให้แพทย์ต้องให้ยาครอบคลุมวงกว้าง (Empirical broad-spectrum therapy) ไปก่อน ซึ่งแนวทางดังกล่าวกลับส่งผลให้ปัญหาเชื้อดื้อยารุนแรงยิ่งขึ้น

วิธี Kirby-Bauer disk diffusion ซึ่งถูกพัฒนาและทดสอบอย่างเป็นระบบตั้งแต่ช่วงทศวรรษ 1960 ยังคงเป็นเทคนิคที่ได้รับการยอมรับอย่างกว้างขวางที่สุดสำหรับการทำ AST ในห้องปฏิบัติการจุลชีววิทยาคลินิกทั่วโลก \cite{asm2009kirby,clsi2025} ความนิยมของวิธีนี้เกิดจากความเรียบง่าย ต้นทุนต่ำ และให้ผลลัพธ์ที่น่าเชื่อถือหากปฏิบัติอย่างถูกต้อง ในขั้นตอนนี้ จานอาหารเลี้ยงเชื้อ Mueller-Hinton agar จะถูกป้ายด้วยสารแขวนลอยของแบคทีเรียเป้าหมาย (มักปรับความขุ่นตามมาตรฐาน 0.5 McFarland) จากนั้นนำแผ่นดิสก์กระดาษสำเร็จรูปที่ชุบด้วยยาปฏิชีวนะในความเข้มข้นที่กำหนด วางลงบนผิววุ้นโดยเว้นระยะห่างตามมาตรฐานเพื่อป้องกันการทับซ้อนกัน หลังจากการบ่มทิ้งไว้ข้ามคืนที่อุณหภูมิ 35--37$^\circ$C ยาปฏิชีวนะจะแพร่กระจายออกจากแผ่นดิสก์เข้าสู่เนื้อวุ้นเป็นรัศมีวงกลม หากความเข้มข้นของยาในวุ้นมีค่าสูงกว่าความเข้มข้นต่ำสุดที่ยับยั้งเชื้อได้ (Minimum Inhibitory Concentration: MIC) การเจริญเติบโตของแบคทีเรียจะถูกยับยั้ง ทำให้เกิดเป็นวงใสรอบแผ่นดิสก์ เรียกว่า บริเวณยับยั้งเชื้อ (Zone of Inhibition: ZOI) ส่วนในบริเวณที่ยาเจือจางจนต่ำกว่าค่า MIC แบคทีเรียก็จะเจริญเติบโตได้ตามปกติ

\begin{figure}[H]
  \centering
  \safeimg[width=0.85\textwidth]{figures/kirby_bauer_example.png}{Example of Kirby-Bauer disk diffusion plate showing zones of inhibition}
  \caption{ตัวอย่างจานเพาะเชื้อแบบ Kirby-Bauer disk diffusion แสดงบริเวณยับยั้งเชื้อ (Zone of Inhibition, ZOI)
    รอบๆ แผ่นดิสก์ยาปฏิชีวนะ เส้นผ่านศูนย์กลางของวงใสแต่ละวงวัดเพื่อ
    จำแนกความไวต่อยาของเชื้อแบคทีเรีย \cite{asm2009kirby}}
  \label{fig:kirby_bauer}
\end{figure}

เส้นผ่านศูนย์กลางของแต่ละ ZOI วัดและนำไปเปรียบเทียบกับจุดตัดการตีความ (interpretive breakpoints) ที่ตีพิมพ์เพื่อจำแนกสิ่งมีชีวิตที่ทดสอบว่าเป็น ไวต่อยา (Susceptible: S), ระดับปานกลาง (Intermediate: I) หรือ ดื้อยา (Resistant: R) จุดตัดเหล่านี้ได้รับการดูแลและปรับปรุงอย่างเข้มงวดโดยองค์กรระหว่างประเทศหลักสองแห่ง ได้แก่ Clinical and Laboratory Standards Institute (CLSI) ในสหรัฐอเมริกา และ European Committee on Antimicrobial Susceptibility Testing (EUCAST) ในยุโรป \cite{eucast2023} CLSI ตีพิมพ์มาตรฐาน M02 ที่ควบคุมขั้นตอนการทดสอบ disk diffusion ในขณะที่ EUCAST ตีพิมพ์ตารางจุดตัดที่แก้ไขเป็นประจำทุกปี ซึ่งได้มาจากแบบจำลอง pharmacokinetic/pharmacodynamic (PK/PD), ค่าจุดตัดทางระบาดวิทยา (epidemiological cutoff values หรือ ECOFFs) และข้อมูลผลลัพธ์ทางคลินิก \cite{giske2022} ทั้งสององค์กรระบุเกณฑ์เส้นผ่านศูนย์กลางเฉพาะสำหรับสายพันธุ์และยาที่ต้องนำไปใช้อย่างถูกต้องเพื่อให้ได้ผลลัพธ์ที่มีความหมายทางคลินิก

แม้จะดูเรียบง่ายในการวัดเส้นผ่านศูนย์กลางของวงกลม แต่ในทางปฏิบัติ การวัด ZOI ด้วยตนเองนั้นมีแนวโน้มที่จะเกิดข้อผิดพลาดได้มาก ช่างเทคนิคในห้องปฏิบัติการใช้ไม้บรรทัด เวอร์เนียร์คาลิปเปอร์ หรืออุปกรณ์เทมเพลตเฉพาะในการวัดแต่ละโซน และความแปรปรวนระหว่างผู้สังเกต (inter-observer variability) เป็นความท้าทายที่ได้รับการยอมรับและยังคงมีอยู่ในสาขานี้ \cite{humphries2018} ปัจจัยต่างๆ เช่น โซนที่ทับซ้อนกันบางส่วนระหว่างแผ่นดิสก์ที่อยู่ติดกัน การเติบโตแบบไม่เป็นเนื้อเดียวกันหรือแบบ ``swarming'' ภายในขอบเขตโซน พื้นผิววุ้นที่ไม่สม่ำเสมอ การปลูกเชื้อบนจานที่ไม่สม่ำเสมอ การฟักตัวที่ไม่เหมาะสม และมุมการมองหรือการถ่ายภาพที่ไม่สอดคล้องกัน ล้วนส่งผลต่อความแปรปรวนในการวัด แม้แต่ความคลาดเคลื่อนเพียงมิลลิเมตรเดียวก็สามารถตัดสินผลทางคลินิกได้ ตัวอย่างเช่น CLSI 2025 กำหนดเส้นผ่านศูนย์กลาง ZOI ที่ 18\,mm หรือมากกว่าสำหรับ \textit{Staphylococcus epidermidis} ที่ทดสอบกับ oxacillin (1\,$\mu$g) ว่าเป็น ไวต่อยา ในขณะที่ 17\,mm หรือน้อยกว่าจะถูกจัดเป็น ดื้อยา เนื่องจากไม่มีหมวดหมู่ Intermediate สำหรับคู่ยา-เชื้อชนิดนี้ ทุกมิลลิเมตรจึงมีความหมาย นอกจากนี้ ห้องปฏิบัติการที่ประมวลผลกรณีจำนวนมากต้องการให้ช่างเทคนิควัดโซนหลายสิบโซนต่อกะการทำงาน ซึ่งทำให้ความเหนื่อยล้ากลายเป็นแหล่งที่มาของข้อผิดพลาดที่เพิ่มมากขึ้น

ระบบ AST อัตโนมัติในปัจจุบันต้องการการลงทุนจำนวนมาก---มักจะเป็นเงินหลายหมื่นดอลลาร์สหรัฐต่อยูนิต---รวมถึงค่าสารเคมีและค่าบำรุงรักษาอย่างต่อเนื่อง ซึ่งเป็นข้อจำกัดสำหรับคลินิกขนาดเล็ก โรงพยาบาลชุมชน และห้องปฏิบัติการในประเทศที่มีรายได้ต่ำและปานกลาง (low- and middle-income countries หรือ LMICs) สภาพแวดล้อมที่มีทรัพยากรจำกัดเหล่านี้มักเป็นที่ที่ได้รับภาระหนักที่สุดจาก AMR ทำให้เกิดความเหลื่อมล้ำอย่างชัดเจนระหว่างห้องปฏิบัติการที่ต้องการ AST ที่แม่นยำมากที่สุดกับห้องปฏิบัติการที่สามารถจ่ายค่าโซลูชันอัตโนมัติได้

ดังนั้น จึงมีเหตุผลที่น่าสนใจในการพัฒนาเครื่องมือซอฟต์แวร์ที่ขับเคลื่อนด้วย AI ราคาประหยัดที่สามารถวิเคราะห์รูปถ่ายมาตรฐานของจาน Kirby-Bauer AST งานวิจัยล่าสุดได้แสดงให้เห็นว่า machine learning และ artificial intelligence สามารถมีบทบาทสำคัญในการทำนายและรับมือกับการดื้อยาต้านจุลชีพ \cite{bilal2025} และความก้าวหน้าในการทดสอบความไวต่อยาที่ขับเคลื่อนด้วย AI กำลังขยายขอบเขตในทางปฏิบัติของ AST อัตโนมัติอย่างรวดเร็ว \cite{liao2025} เครื่องมือดังกล่าวสามารถติดตั้งบนฮาร์ดแวร์ทั่วไป เช่น แล็ปท็อปหรือแม้แต่สมาร์ทโฟน และจะไม่ต้องการเครื่องมือเฉพาะทางนอกเหนือจากจานวุ้นเอง ซึ่งอาจช่วยขยายขอบเขตของ AST ที่แม่นยำไปยังสภาพแวดล้อมที่ปัจจุบันต้องพึ่งพาการอ่านด้วยตนเอง รวมถึงสภาพแวดล้อมที่ได้รับคำแนะนำจากมาตรฐานห้องปฏิบัติการระดับชาติ เช่น มาตรฐานที่ดูแลโดยกรมวิทยาศาสตร์การแพทย์ สถาบันวิจัยวิทยาศาสตร์สาธารณสุข (NIH Thailand) \cite{dmsc_nih} โครงงานปัจจุบันนี้ตอบสนองความต้องการนี้โดยการพัฒนา ประเมิน และติดตั้งระบบที่ใช้ AI สำหรับการตรวจจับและวัด ZOI อัตโนมัติ

\section{แรงจูงใจ (Motivation)}

ความก้าวหน้าล่าสุดในสาขา Deep Learning และ Computer Vision ได้ก่อให้เกิดโมเดลที่มีประสิทธิภาพเทียบเท่าหรือเหนือกว่ามนุษย์ในงานด้านการจดจำรูปภาพ (Image Recognition) และการตรวจจับวัตถุ (Object Detection) ที่หลากหลาย สำหรับการวิเคราะห์ภาพ AST นั้น ความก้าวหน้าที่เกี่ยวข้องโดยตรงคือความสามารถในการตรวจจับและระบุตำแหน่งของโครงสร้างวงกลมหรือวงรีบนพื้นหลังที่ซับซ้อนและมีสัญญาณรบกวน ต่างจากชุดข้อมูลภาพถ่ายทั่วไปที่วัตถุมีรูปร่างหลากหลาย ภาพถ่ายจานเพาะเชื้อ AST มีลักษณะทางกายภาพที่ค่อนข้างคงที่ นั่นคือ ประกอบด้วยโซนวงกลมบนพื้นหลังวุ้นที่มีพื้นผิวเฉพาะ พร้อมด้วยวัตถุอ้างอิงที่มีขนาดตายตัว (แผ่นดิสก์ยาปฏิชีวนะ) ปรากฏอยู่ในทุกภาพ โครงงานนี้จึงดึงเอาข้อได้เปรียบทางโครงสร้างภาพดังกล่าวมาใช้ในการออกแบบระบบ

แรงจูงใจหลักของโครงงานนี้เกิดจากปัจจัยสองประการ ประการแรก โมเดล Object Detection ที่ใช้เทคนิค Deep Learning---โดยเฉพาะสถาปัตยกรรมแบบ Two-stage อย่าง Faster R-CNN และแบบ Single-stage อย่าง YOLOv11---ได้พิสูจน์ให้เห็นถึงความสามารถในการเรียนรู้เพื่อระบุ ZOI จากชุดข้อมูลภาพถ่ายได้อย่างแม่นยำ แม้ในกรณีที่โซนมีการทับซ้อนกันหรือมีรูปร่างผิดปกติ ประการที่สอง โซลูชันซอฟต์แวร์สำหรับการวิเคราะห์ AST ที่สมบูรณ์แบบนั้น ไม่เพียงแต่ต้องตรวจจับโซนได้เท่านั้น แต่ยังต้องสามารถเชื่อมโยงแต่ละโซนเข้ากับแผ่นดิสก์ยาปฏิชีวนะที่ถูกต้องได้ด้วย หากไม่ทราบว่ายาปฏิชีวนะชนิดใดเป็นตัวสร้างโซนนั้น การวัดขนาดเส้นผ่านศูนย์กลางย่อมไม่มีความหมายทางคลินิก การแก้ปัญหาการเชื่อมโยงนี้จำเป็นต้องอาศัยการอ่านรหัสตัวอักษรและตัวเลขที่พิมพ์อยู่บนแผ่นดิสก์ ซึ่งเป็นงานหลักของเทคโนโลยีการจดจำอักขระด้วยแสง (Optical Character Recognition: OCR)

การบูรณาการ Object Detection Model สำหรับระบุตำแหน่ง ZOI, ระบบ OCR สำหรับอ่านชื่อยาปฏิชีวนะ และฐานข้อมูลเกณฑ์จุดตัด (Breakpoints) สำหรับการจำแนกผลทางคลินิก ถือเป็นแนวคิดหลักในการพัฒนาระบบนี้ องค์ประกอบแต่ละส่วนล้วนมีความท้าทายทางเทคนิคที่แตกต่างกันไป การตรวจจับวัตถุต้องจัดการกับสัญญาณรบกวน แสงที่แปรปรวน และการทับซ้อนของโซน ในขณะที่ OCR ต้องอ่านตัวอักษรขนาดเล็กที่มีความเปรียบต่างต่ำบนแผ่นดิสก์ที่อาจโค้งงอและถูกถ่ายภาพจากมุมใดก็ได้ ส่วนระบบการปรับเทียบ (Calibration) ต้องสามารถแปลงหน่วยพิกเซลให้เป็นหน่วยมิลลิเมตรได้อย่างถูกต้อง โดยอ้างอิงจากขนาดจริงของแผ่นดิสก์ที่ปรากฏในภาพ

นอกจากนี้ โครงงานนี้ยังมุ่งเปรียบเทียบประสิทธิภาพระหว่างเทคนิคประมวลผลภาพแบบคลาสสิกและวิธีการเรียนรู้เชิงลึก (Deep Learning) สมัยใหม่ ภายใต้ชุดข้อมูลและเกณฑ์การประเมินเดียวกัน โดยใช้วิธี Hough Circle Transform ซึ่งเป็นเทคนิคดั้งเดิมสำหรับการตรวจจับวงกลม เป็นเกณฑ์มาตรฐาน (Baseline) เพื่อแสดงให้เห็นถึงประโยชน์และข้อได้เปรียบของ Deep Learning สำหรับงานวิเคราะห์ภาพลักษณะนี้อย่างเป็นรูปธรรม

\section{วัตถุประสงค์ (Objectives)}

โครงงานนี้มีวัตถุประสงค์หลักสี่ประการ:
\begin{enumerate}
  \item เพื่อศึกษาและดำเนินการตีความผลการทดสอบ Kirby-Bauer disk diffusion ตามมาตรฐาน CLSI 2025 \cite{clsi2025} และ EUCAST 2023 \cite{eucast2023} เพื่อให้มั่นใจว่าการจำแนกความไวต่อยาของระบบมีพื้นฐานทางคลินิกที่มั่นคง
  \item เพื่อพัฒนาและฝึกฝนโมเดล AI---โดยเฉพาะ Faster R-CNN และ YOLOv11n---ที่สามารถตรวจจับ ZOI และวัดเส้นผ่านศูนย์กลางจากภาพถ่ายจาน AST ด้วยค่าความคลาดเคลื่อนสมบูรณ์เฉลี่ย (mean absolute error หรือ MAE) ที่หรือต่ำกว่าเกณฑ์ที่ยอมรับได้ทางคลินิกที่ 3\,mm \cite{humphries2018}
  \item เพื่อเปรียบเทียบความแม่นยำ ความเร็ว และความทนทานของแนวทางการตรวจจับสามแบบอย่างเข้มงวด: Faster R-CNN (deep learning แบบสองขั้นตอน), YOLOv11n (deep learning แบบขั้นตอนเดียว) และ Hough Circle Transform (เกณฑ์มาตรฐานการประมวลผลภาพแบบคลาสสิก) เพื่อเป็นฐานเชิงปริมาณในการเลือกโมเดลที่มีประสิทธิภาพดีที่สุดสำหรับการติดตั้งใช้งาน
  \item เพื่อติดตั้งใช้งานระบบเป็นเว็บแอปพลิเคชันที่ใช้งานได้จริงและเป็นมิตรต่อผู้ใช้ ซึ่งเหมาะสมสำหรับการใช้งานจริงในสภาพแวดล้อมห้องปฏิบัติการทางคลินิก โดยผสานรวมโมเดลการตรวจจับเข้ากับไปป์ไลน์ OCR และฐานข้อมูลการจำแนกจุดตัดให้เป็นเครื่องมือที่สอดคล้องและเข้าถึงได้
\end{enumerate}

\section{ขอบเขตของโครงงาน (Project Scope)}

โครงงานนี้มีขอบเขตการดำเนินงาน ดังนี้:

\begin{itemize}
  \item \textbf{ข้อมูลนำเข้า (Input Data):} ระบบรองรับภาพถ่ายจานเพาะเชื้อ Mueller-Hinton ที่ใช้ทดสอบ AST แบบ Kirby-Bauer ตามมาตรฐาน โดยภาพควรถูกถ่ายจากมุมมองด้านบน (Top-down) มีแสงสว่างเพียงพอ และระยะห่างระหว่างกล้องกับจานเพาะเชื้อค่อนข้างคงที่ ภาพสามารถถ่ายจากกล้องในห้องปฏิบัติการหรือกล้องสมาร์ทโฟนทั่วไปได้

  \item \textbf{แอปพลิเคชันที่พร้อมใช้งาน (Deployed Application):} แอปพลิเคชันที่ทำงานได้อย่างสมบูรณ์ซึ่งรวมโมเดลการตรวจจับที่มีประสิทธิภาพดีที่สุด Pipeline การจดจำชื่อยาปฏิชีวนะที่ใช้ EasyOCR และฐานข้อมูลการแปลผลตามเกณฑ์จุดตัด CLSI 2025/EUCAST 2023 มอบเวิร์กโฟลว์แบบครบวงจรตั้งแต่การอัปโหลดภาพไปจนถึงรายงานความไวต่อยา

  \item \textbf{ชุดข้อมูลที่ผ่านการ Label (Annotated Dataset):} ชุดข้อมูลที่ผ่านการคัดสรรและ Label (Annotated) ของภาพจาน AST จริงจำนวน 287 ภาพ ซึ่งขยายเป็น 543 ภาพผ่านเทคนิค Augmentation และเสริมด้วยภาพสังเคราะห์ 2,000 ภาพ เป็นทรัพยากรที่มีค่าสำหรับการวิจัย AI ในอนาคต


  \item \textbf{การปรับเทียบ (Calibration):} การแปลงหน่วยจากพิกเซลเป็นมิลลิเมตรจะทำโดยใช้แผ่นดิสก์ยาปฏิชีวนะในภาพเป็นวัตถุอ้างอิง เนื่องจากแผ่นดิสก์มาตรฐานมีเส้นผ่านศูนย์กลางที่กำหนดไว้ชัดเจนคือ 6\,mm ตามข้อกำหนดของ CLSI ทำให้ไม่จำเป็นต้องใช้ไม้บรรทัดหรือวัตถุอ้างอิงภายนอกอื่นในการถ่ายภาพ

  \item \textbf{เว็บแอปพลิเคชัน (Web Application):} ผู้ใช้สามารถอัปโหลดภาพหรือถ่ายภาพแบบเรียลไทม์เข้าสู่ระบบได้ ระบบจะแสดงผลในรูปแบบ Dashboard ประกอบด้วย ZOI ที่ตรวจพบ, ชื่อยาปฏิชีวนะที่อ่านได้จาก OCR, ขนาดเส้นผ่านศูนย์กลาง, และผลการจำแนก (S/I/R) ตามมาตรฐาน CLSI หรือ EUCAST พร้อมระบบบันทึกประวัติเพื่อตรวจสอบย้อนหลัง

  \item \textbf{ข้อจำกัดของระบบ (Out of Scope):} ระบบไม่รองรับการวิเคราะห์หรือระบุสายพันธุ์ของแบคทีเรียอัตโนมัติ ผู้ใช้งานต้องเป็นผู้กำหนดชนิดของเชื้อ (เช่น \textit{Staphylococcus aureus}, \textit{Escherichia coli}) ด้วยตนเองเพื่อให้ระบบเลือกใช้ตารางเกณฑ์แปลผลได้อย่างถูกต้อง นอกจากนี้ ระบบไม่รองรับการทดสอบแบบ Broth Microdilution, E-test หรือรูปแบบอื่นๆ และผลลัพธ์จากระบบมีจุดประสงค์เพื่อสนับสนุนการตัดสินใจเท่านั้น ไม่ใช่เพื่อทดแทนการวินิจฉัยโดยบุคลากรทางการแพทย์
\end{itemize}

\section{ผลที่คาดว่าจะได้รับ (Expected Outcomes)}

โครงงานนี้คาดว่าจะสร้างผลลัพธ์ที่เป็นรูปธรรมดังนี้:

\begin{enumerate}
  \item \textbf{ไปป์ไลน์การตรวจจับด้วย AI (AI Detection Pipeline):} ได้ระบบประมวลผลอัตโนมัติแบบครบวงจรที่สามารถระบุตำแหน่งและวัดขนาดเส้นผ่านศูนย์กลาง ZOI จากภาพถ่าย AST โดยโมเดลที่ดีที่สุดจะต้องมีค่า MAE ต่ำกว่า 3\,mm บนชุดข้อมูลทดสอบ ซึ่งอยู่ในเกณฑ์ที่ยอมรับได้ทางคลินิก \cite{humphries2018}

  \item \textbf{การวิเคราะห์เปรียบเทียบโมเดล (Comparative Model Analysis):} ได้ผลการเปรียบเทียบเชิงปริมาณอย่างเป็นระบบระหว่าง Faster R-CNN, YOLOv11n และ Hough Circle Transform เพื่อนำเสนอจุดเด่นและจุดด้อยของแต่ละวิธี ทั้งในด้านความแม่นยำ ความเร็ว และความซับซ้อน เพื่อเป็นแนวทางสำหรับงานวิจัยด้านวิเคราะห์ภาพทางการแพทย์ในอนาคต

  \item \textbf{เว็บแอปพลิเคชันพร้อมใช้งาน (Deployed Web Application):} ได้แอปพลิเคชันที่ทำงานได้อย่างสมบูรณ์ โดยบูรณาการทั้งโมเดล AI, ระบบ OCR ด้วย EasyOCR และฐานข้อมูลมาตรฐาน CLSI 2025/EUCAST 2023 เข้าด้วยกันเพื่อให้บริการการวิเคราะห์ตั้งแต่ต้นจนจบ (End-to-end Workflow)

  \item \textbf{ชุดข้อมูลพร้อมป้ายกำกับ (Annotated Dataset):} ได้ชุดข้อมูลภาพถ่ายจาน AST จริงพร้อมป้ายกำกับ (Ground-truth) จำนวน 287 ภาพ รวมถึงภาพที่ผ่านการขยายข้อมูล (Augmented) จำนวน 543 ภาพ และภาพสังเคราะห์ (Synthetic) อีก 2,000 ภาพ ซึ่งจะเป็นทรัพยากรที่มีค่าสำหรับการวิจัยและพัฒนา AI ในโดเมนนี้ต่อไป
\end{enumerate}

\section{โครงสร้างของรายงาน (Report Structure)}

ส่วนที่เหลือของรายงานนี้จัดลำดับดังนี้: บทที่ 2 ทบทวนทฤษฎีพื้นฐานและงานวิจัยที่เกี่ยวข้อง ครอบคลุมวิธี Kirby-Bauer AST และมาตรฐานทางคลินิก สถาปัตยกรรมและหลักการของโมเดลการตรวจจับวัตถุ (Faster R-CNN, YOLOv11n) และ Hough Circle Transform, ระบบ OCR สำหรับระบุชื่อยาปฏิชีวนะ, เทคนิคการประมวลผลภาพเบื้องต้นที่สำคัญ และการทบทวนงานวิจัยเกี่ยวกับการวิเคราะห์ภาพ AST อัตโนมัติ บทที่ 3 อธิบายระเบียบวิธีวิจัยอย่างละเอียด รวมถึงขั้นตอนการรวบรวมและเตรียมชุดข้อมูล การ Label ข้อมูล (Annotation) กระบวนการฝึกโมเดล แนวทางการปรับเทียบ การออกแบบ Pipeline OCR ตรรกะการจำแนกความไวต่อยา และสถาปัตยกรรมระบบโดยรวม บทที่ 4 นำเสนอผลการทดลอง รวมถึงความแม่นยำในการตรวจจับ การทดสอบความเร็ว ประสิทธิภาพของ OCR และการวิเคราะห์ปัญหาในการปรับเทียบขนาด ท้ายสุด บทที่ 5 สรุปสิ่งที่ค้นพบ อภิปรายข้อจำกัด และเสนอแนวทางสำหรับงานในอนาคต
